\section*{Chapitre \ref{ch:s2:capital-structure-arbitrage} --- Arbitrage de structure du capital}

\begin{solution}{exo:s2:capital-structure-arbitrage:1}
$0.6 \times (-\ln 0.954)/5 = 0.00565$~: 56.5 points de base.
\end{solution}

\begin{solution}{exo:s2:capital-structure-arbitrage:2}
$0.00067 \times 10\,000\,000 = 6\,700$ actions, soit environ 0.34~million de dollars au prix de 50.
\end{solution}

\begin{solution}{exo:s2:capital-structure-arbitrage:3}
$4.5 \times 0.0043 = 0.0194$~: 193.5 points de base de notionnel, avant le carry perçu.
\end{solution}

\begin{solution}{exo:s2:capital-structure-arbitrage:4}
De nombreux écarts proviennent de la mauvaise barrière retenue par le trader, et non d'un défaut de prix~: la vraie valeur de $L$ de l'entreprise diffère de 0.5 supposé, de sorte que le spread de marché reste
durablement éloigné du \omterm{def:s2:capital-structure-arbitrage:implied}{spread implicite actions}, et l'écart ne se referme jamais.
\end{solution}

\begin{solution}{exo:s2:capital-structure-arbitrage:5}
La dette double, si bien que le vrai spread bondit et que le vendeur de protection perd sur le CDS~; l'action bondit sous l'effet de la prime de rachat, de sorte que la couverture à découvert
perd elle aussi. Le modèle supposait que les actions et le crédit évoluent en sens contraire~; un rachat les fait évoluer dans le même sens.
\end{solution}

\begin{solution}{exo:s2:capital-structure-arbitrage:6}
Les spreads et les cours des actions sont faiblement corrélés au jour le jour, de sorte que le P\&L de chaque transaction est bruité et peut être important~; sur des centaines de transactions, le bruit
s'annule en moyenne et la correction de l'écart de prix subsiste.
\end{solution}

\begin{solution}{exo:s2:capital-structure-arbitrage:7}
La pire transaction passe de $-347$ à $-179$ points de base et le ratio de Sharpe passe de 1.40 à 1.66~; l'erreur de modèle sur la barrière subsiste, et
44~\% des transactions restent perdantes.
\end{solution}

\begin{solution}{exo:s2:capital-structure-arbitrage:8}
Un écart qui persiste pendant six mois résulte plus probablement d'une mauvaise donnée d'entrée (la barrière, la dette, la volatilité) ou d'un changement dans l'entreprise, plutôt
qu'un défaut de prix en attente de correction~; dans la simulation du chapitre, les écarts qui ne convergent pas sont des erreurs du modèle.
\end{solution}

\begin{solution}{pb:s2:capital-structure-arbitrage:1}
\begin{enumerate}
\item Actifs $E + LD$, volatilité des actifs $\sigma_E E/V$, défaut lorsque les actifs touchent $LD$~; données~: cours de l'action, volatilité des actions, dette par action, $L$ et taux de recouvrement.
\item Le spread de CDS que le modèle déduit des actions~: 55.8 points de base avec une dette de 40, et 115.4 avec une dette de 80.
\item La volatilité des actifs du modèle diminue lorsque le cours de l'action augmente.
\item CreditGrades, 135\,759 spreads quotidiens sur 261 entreprises~; pertes importantes sur une transaction isolée, ratio de Sharpe de portefeuille comparable à celui des autres stratégies d'arbitrage à revenu fixe.
\item Nombre d'actions par unité de notionnel de CDS compensant le mouvement du CDS par rapport au cours de l'action~; environ 6\,700 actions, soit 0.34~million de dollars.
\item Entrer à un écart de 40~\% en logarithme~; sortir sous 5~\% ou après 180 jours.
\item Sa vraie barrière et, pendant trois mois après un rachat, sa dette.
\item La jambe CDS~: 20.4 des 26.1 points de base par transaction.
\item 440 transactions, 26.1 points de base chacune, 48~\% perdantes, pire résultat $-347$, ratio de Sharpe 1.40.
\item Leurs écarts relevaient d'une erreur de modèle.
\item Les stratégies qui exigeaient davantage de capital intellectuel ont dégagé des alphas significatifs, avec des rendements à asymétrie positive.
\item Ratio de Sharpe 1.66, pire transaction $-179$.
\item \textbf{Résultat clé.} Le portefeuille de convergence dégage un ratio de Sharpe de 1.40, avec 48~\% des transactions perdantes~; les vendeurs de protection pris dans un
rachat avec effet de levier ont perdu 98 points de base en moyenne, et avec des rachats quatre fois plus fréquents, le ratio de Sharpe tombe à 0.39.
\item Les vendeurs de protection ont perdu 98 points de base en moyenne, les acheteurs ont gagné 220.
\item Elle n'est pas observable et varie d'une entreprise à l'autre~; la stratégie la traite comme connue.
\item Fermer les écarts qui ne convergent pas, écarter les entreprises faisant l'objet d'une offre, mettre à jour la dette à partir des dépôts réglementaires, plafonner la taille des positions.
\item Dette et volatilité connues à l'époque, coûts, transactions clôturées au cours limite, rachats présents dans l'échantillon.
\item Désalignement dette--actions après des événements.
\item Il reprend la probabilité de survie du modèle de premier passage du livre~6 pour convertir les actions en spread.
\item Que les actions et le crédit s'accorderont de nouveau sur le risque de défaut de l'entreprise.
\end{enumerate}
\end{solution}

\begin{solution}{iq:s2:capital-structure-arbitrage:1}
Avec un modèle structurel~: traiter les actions comme une créance sur les actifs, en déduire la valeur et la volatilité des actifs à partir du cours de l'action et de la volatilité des actions, fixer une barrière de défaut
à partir de la dette, calculer la probabilité de survie et la convertir en spread.
\end{solution}

\begin{solution}{iq:s2:capital-structure-arbitrage:2}
Vendre une protection et vendre à découvert des actions au ratio de couverture du modèle~; encaisser le carry et la convergence. Risques~: le modèle est faux (barrière, volatilité, dette), un
rachat ou une recapitalisation fait évoluer les actions et le crédit dans le même sens, le prêt de titres fait l'objet d'un rappel, et l'écart s'élargit avant de se refermer.
\end{solution}

\begin{solution}{iq:s2:capital-structure-arbitrage:3}
Un changement de l'effet de levier sur les plus grosses positions, un élargissement des spreads à l'échelle du marché alors que les actions restent stables, une variation de la volatilité des actions, et un rappel de prêt de titres~; avec
le notionnel brut du portefeuille et la perte sur les transactions qui ne convergent jamais.
\end{solution}

\begin{solution}{iq:s2:capital-structure-arbitrage:4}
Un défaut de prix se corrige~; une erreur de modèle persiste ou suit une donnée d'entrée. Il faut vérifier si l'écart évolue avec l'actualité propre à l'entreprise ou ses dépôts réglementaires, tester d'autres
barrières, et examiner les obligations et options émises par la même entreprise.
\end{solution}

\begin{solution}{iq:s2:capital-structure-arbitrage:5}
Cotations de CDS, cours des actions, volatilité des actions et prêt de titres au quotidien~; dette par action issue des dépôts trimestriels (publiés avec retard)~; événements d'entreprise tels qu'annoncés. La dette est la
donnée d'entrée la plus susceptible d'être périmée.
\end{solution}

\begin{solution}{iq:s2:capital-structure-arbitrage:6}
Pour $X_t = \ln V_t$ avec dérive $\nu = -\sigma^2/2$ et $x = \ln(V_0/B)$, le principe de réflexion pour le mouvement brownien avec dérive donne
$P(\min_{u \le t} X_u > \ln B) = N\big((x + \nu t)/(\sigma\sqrt t)\big) - e^{-2\nu x/\sigma^2}N\big((-x + \nu t)/(\sigma\sqrt t)\big)$, et avec
$\nu = -\sigma^2/2$, le facteur $e^{-2\nu x/\sigma^2}$ vaut $e^{x} = V_0/B$.
\end{solution}
